\begin{pastexam}[label={px:ust-114-ee-ds-16}]{臺聯大電機類 114 學年度 第 16 題}
（單選；2 分，答錯不倒扣）

Answer questions 14-16 according to the following C code.
\begin{minted}[fontsize=\small]{c}
struct Node {
    int data;
    struct Node* next;
};
int funcA(struct Node* head) {
    struct Node *slow = head, *fast = head;
    while (slow && fast && fast->next) {
        slow = slow->next;
        fast = fast->next->next;
        if (slow == fast) {
            return 1;
        }
    }
    return 0;
}
\end{minted}

Auxiliary memory space required by funcA()?

(A)~O(1) \qquad (B)~O(n) \qquad (C)~O(n$^2$) \qquad (D)~O(nlogn)
\end{pastexam}

\begin{pastsoln}
答案是 (A)。

\texttt{funcA} 只用了兩個指標變數 \texttt{slow}、\texttt{fast}，沒有遞迴、沒有配置額外的陣列或雜湊
表，額外空間和 $n$ 無關，是 $O(1)$。 這正是 Floyd 判環法的優點：若改用「把走過的節點記在雜湊集合
裡，看有沒有重複」的做法，時間一樣 $O(n)$，空間卻要 $O(n)$。 (B)(C)(D) 都錯。
\end{pastsoln}
